\documentclass[10pt,a4paper]{amsart}
\usepackage[margin=19mm]{geometry}
\usepackage{amsmath,amssymb,mathtools}
\usepackage[scr=rsfs,cal=euler]{mathalfa}
\usepackage{xfrac}
\usepackage[unicode,hidelinks,bookmarks=false]{hyperref}
\newcommand{\ie}{i.e.\ }
\newcommand{\cf}{cf.\ }
\newcommand{\resp}{respectively\ }
\newcommand{\Subsection}{\subsection}
\providecommand{\nobreakdash}{\nobreak\hbox{-}}
\setlength{\emergencystretch}{2em}
\multlinegap0pt
\usepackage{amssymb}
\usepackage{mathtools}
\usepackage{float}
\newtheorem{proposition}{Proposition}
\newtheorem{lemma}{Lemma}
\newcommand{\R}{\mathbb{R}} 
\title{Solving the Offline and Online Min-Max Problem of Non-smooth Submodular-Concave Functions: A Zeroth-Order Approach}
\author{Amir Ali Farzin, Yuen-Man Pun, Philipp Braun, Tyler Summers, Iman Shames}
\date{Source: arXiv:2601.21243v4 (CC BY 4.0)}
\begin{document}
\maketitle

\noindent\emph{Source and licence.} Solving the Offline and Online Min-Max Problem of Non-smooth Submodular-Concave Functions: A Zeroth-Order Approach. Version arXiv:2601.21243v4 (CC BY 4.0), distributed by arXiv under the Creative Commons Attribution 4.0 International licence, http://creativecommons.org/licenses/by/4.0/, as recorded in the arXiv licence metadata for this exact version and shown on the abstract page https://arxiv.org/abs/2601.21243v4. Full licence notice: \texttt{licenses/arxiv-2601.21243-CC-BY-4.0.txt}.\par\smallskip

\noindent\textit{Editorial scope.} Counted unit: the article's proposition Prp:sadex (source0.tex lines 394--412). The article's complete original proof of it is delivered with it (source0.tex lines 1045--1048, one \texttt{proof} environment of 4 lines, which the article places away from the statement). No mathematics is written, repaired or re-derived: the statement and the proof are the article's own lines, in the article's own notation, and no retained line is substituted or re-flowed.  Retained uncounted as the source-local context the proof needs: lines 388--390; lines 892--894; lines 1035--1037.  Nothing was omitted from any retained span, so the package carries no omission record.  No retained line contains a citation, so the wrapper carries no bibliography and no bibliography entry.  No retained line contains a figure environment, an image-inclusion command or drawing code, so no image is delivered and the package declares no asset dependency.  Editorial formatting only, in addition: an AMS \texttt{amsart} preamble that restates the article's own declarations and macros used by the retained lines, namely the environments \texttt{proposition}, \texttt{lemma} on separate counters, together with the standard packages those lines need.  The retained environments print in this document as Proposition 1, Lemma 1 in order of appearance, whereas the article prints the same items with the numbering of its own full document.  The complete native source of the article as distributed in the arXiv e-print for this exact version is bundled unchanged as \texttt{sources/193578/source0.tex}.
\begin{align}\label{eq:ver1}
    \underset{x \in [0,1]^n}{\min}\underset{y \in \mathcal{Y}}{\max}\;f^L(x,y)=\underset{x \in \{0,1\}^n}{\min}\underset{y \in \mathcal{Y}}{\max}\;f^L(x,y).
\end{align}

\begin{lemma}\label{lm:lmlova}
    Let $\bar{f}:2^{[n]}\rightarrow \mathbb{R}$ be a submodular function and $\bar{f}^L$ be its Lovász extension. Then, we have that (i) $\bar{f}^L$ is convex, (ii) $\bar{f}^L$ is piecewise linear,  (iii) $\bar{f}^L$ is positively homogeneous (i.e., for each $c>0,$ we have $\bar{f}^L(cx)=c\bar{f}^L(x)$),  (iv) $\min_{S\subset[n]} \bar{f}(x) = \min_{x\in\{0,1\}^n} \bar{f}^L(x)=\min_{x\in[0,1]^n} \bar{f}^L(x)$, and (v) the set of minimisers of $\bar{f}^L(x)$ in $[0,1]^n$ is the convex hull of the set of minimisers of $\bar{f}^L(x)$ in $\{0,1\}^n$. 
\end{lemma}

\begin{align}\label{eq:ver}
    \underset{x \in [0,1]^n}{\min}\underset{y \in \mathcal{Y}}{\max}\;f^L(x,y)=\underset{x \in \{0,1\}^n}{\min}\underset{y \in \mathcal{Y}}{\max}\;f^L(x,y)
\end{align}

\begin{proposition}\label{Prp:sadex}
    Let $f:2^{[n]}\times\R^m\rightarrow \mathbb{R}$ be a submodular-concave function and $f^L$ be the Lovász extension of $f$ with respect to the first variable. Let 
    %$\zeta(x) = \underset{y \in \mathcal{Y}}{\max}\;f^L(x,y)$
    $\zeta(x) = \max_{y \in \mathcal{Y}}\;f^L(x,y)$ 
    for $\mathcal{Y}\subset \R^m$ convex and compact. Then, \eqref{eq:ver1} %\eqref{eq:ver} 
    holds, and the existence of a saddle point for $f$ is guaranteed if one of the following conditions is satisfied: i)  $\zeta(x)$ itself is a Lovász extension of a submodular function. ii) The set of minimisers of $\zeta(x)$ contains a vertex. iii) The family of functions $\{f^L(\cdot,y):y\in\mathcal{Y}\}$ has a common
    vertex minimiser.
    % % which is a convex function as $f^L$ is convex with respect to $x.$ 
    % \begin{itemize}
    %     \item[i)] 
    % %If 
    %  $\zeta(x)$ itself is a Lovász extension of a submodular function. %or 
    % \item [ii)] The set of minimiser of $\zeta(x)$ contains a vertex. %or 
    % % iii) there exists a vertex $\chi_{S^\ast}$ that simultaneously minimises every $f^L(\cdot,y)$ that is active at the optimum of the max, or 
    % \item[iii)] The family of functions $\{f^L(\cdot,y):y\in\mathcal{Y}\}$ has a common
    % vertex minimiser.
    % \end{itemize}
    % %then \eqref{eq:ver} holds and existence of a saddle point for $f$ is guaranteed.
\end{proposition}

\begin{proof}[Proof of Proposition~\ref{Prp:sadex}]
    To prove i), considering Lemma~\ref{lm:lmlova}.v), if $\zeta$ is a Lovász extension of a submodular function, then \eqref{eq:ver} holds immediately. Moreover, if the set of minimiser of $\zeta(x)$ over $x\in[0,1]^n$ contains a vertex $x\in\{0,1\}^n,$ it is trivial to see that \eqref{eq:ver} holds.
    To prove the last statement, if for any $y\in\mathcal{Y},$ the set of minimiser of $f(\cdot ,y)$ share a common vertex $x\in\{0,1\}^n,$ then it is guaranteed that the set of minimiser of $\zeta(x)$ contains a vertex and, consequently, \eqref{eq:ver} holds.
\end{proof}

\end{document}
